% Primitive: cantor-diagonal — authored from
% probability/cantor_diagonal_manim.py, its README concepts table (seven
% scenes) and plan 017's verified anchors (verifier report A–K plus the
% addendum L–T). Follows Bertsekas & Tsitsiklis's Problem 4* in its own
% language: "uncountable" means "cannot be arranged in a sequence", and
% nothing heavier is claimed. Problem answers are computed by
% answers/cantor_diagonal.py (the solve-gate anchor); every digit there is
% integer long division, never a float.
\chapter{Cantor's diagonal}

\begin{narrative}
Every chapter so far has counted something and then reweighted the
count into area. This chapter is about the one set the repo's pictures
have quietly relied on and never counted: the interval $[0,1]$, the
sample space of the wheel of fortune. Bertsekas and Tsitsiklis pose it
as a starred problem — \emph{show that the unit interval $[0,1]$ is
uncountable, i.e., its elements cannot be arranged in a
sequence}~\cite{probability-bertsekas-and-tsitsiklis-introduction-to-probability-2nd-ed} —
and the whole chapter is their solution, slowed down until each step is
a picture. The payoff lands back in probability: the additivity axiom is
stated for \emph{sequences} of events, and the interval is not one.

\section{Arranged in a sequence}

``Arranged in a sequence'' means exactly this: every element gets a
finite position number. Six tokens take positions one to six, the row
keeps going, and nothing stops a point from being named twice —
repeats are allowed. Velleman calls the property
\emph{listability}~\cite{probability-daniel-j-velleman-how-to-prove-it-2nd-ed-cambridge-2006-ch-7},
and the surprise of the chapter is how much can be listed. Stack the
integers under the positions and pair them off: even $n$ goes to $n/2$,
odd $n$ to $-(n-1)/2$, so positions $1, 2, 3, \ldots, 10$ hold
$\anchor{017.P.pairs}$ and every integer has exactly one position. The
lists are drawn stacked, one above the other, with an arrow per column —
never side by side — because side-by-side sets invite the part–whole
reading (``the integers are bigger, they contain the naturals'') that
the pairing is meant to
replace~\cite{probability-tsamir-and-tirosh-consistency-and-representations-jrme-30-19}.
One more fact, stated now because the diagonal will need it later: a
list can always take one more. Shift every element one place to the
right and a new element fits at position one — Hilbert's
hotel~\cite{probability-wikipedia-hilbert-s-paradox-of-the-grand-hotel}.
Running out of room is never the obstacle.

\section{The zigzag}

Pairs of positions $(p, q)$ are the counting grid of the first chapter,
now open on two sides. Walk it along the anti-diagonals — $p + q = 2$,
then $3$, then $4$, starting each diagonal at the top row — and stamp
the cells $1, 2, 3, \ldots$ as you go: the first row reads $1, 2, 4, 7,
11$, and the $\anchor{017.O.fifteen}$ cells with $p + q \le 6$ are
stamped $1$ to $15$. Every cell gets a finite position, so pairs are a
sequence. Now read each cell as the fraction $p/q$. Some fractions
appear twice — $2/2$ is $1/1$ again, $2/4$ is $1/2$ — so skip any cell
not in lowest terms. Among the first fifteen stamps four are skipped
and $\anchor{017.O.kept}$ are kept, and the kept ones take positions
$1$ to $11$ in walk order: $1/1, 1/2, 2/1, 1/3, 3/1, 1/4, 2/3, 3/2,
4/1, 1/5, 5/1$. Skipping never breaks a list: the next kept cell takes
the next position. Every positive fraction has exactly one lowest-terms
cell, so exactly one position — the rationals are a
sequence~\cite{probability-wikipedia-countable-set}. This is the beat
that earns the surprise to come: ``more than the naturals'' in every
naive sense is still a sequence.

\section{Every point has digits}

Cut the unit segment into tenths. The point $1/3$ lies in the fourth
tenth, $[0.3, 0.4)$, so its first digit is $3$; cut that tenth into
tenths again and it lies in the fourth of those, so the second digit is
$3$ — and so on for ever, $1/3 = 0.333\cdots$, the example the source
itself uses. The digits are the point's address, and every point has
one~\cite{probability-michigan-state-math-320-decimal-expansions-of-real-numbers-2}.
A few points have two. The point $1/2$ sits on a tick, the shared edge
of $[0.4, 0.5)$ and $[0.5, 0.6)$, and it can be written either way:
$0.5000\cdots$ or $0.4999\cdots$. That these are the same number is
$1/3$'s example again — three times $0.333\cdots$ is $0.999\cdots$, and
three times $1/3$ is $1$. Bertsekas and Tsitsiklis state, without
proof, that ``this is the only kind of exception'': the rows that end
in an infinite string of zeros or of
nines~\cite{probability-stevecheng-existence-and-uniqueness-of-decimal-expansion-pla}.
Hold on to the exception; the next section but one is about it.

\section{The diagonal rule}

Suppose the points of $[0,1]$ \emph{are} arranged in a sequence $x_1,
x_2, x_3, \ldots$, and write each as its row of digits, in the source's
notation $x_n = 0.a_n^1 a_n^2 a_n^3 \cdots$ — the subscript says which
number, the superscript which digit. Seven rows stand in for the
sequence: $\pi - 3 = \anchor{017.L.pi}\ldots$, then $1/2$ as
$0.5000\ldots$, $1/3$, $\sqrt2 - 1 = \anchor{017.L.sqrt}\ldots$, $1/e =
\anchor{017.L.inve}\ldots$, $e - 2 = \anchor{017.L.eminus}\ldots$, and
$1/2$ once more, deliberately, under its other name $0.4999\ldots$.
Read digit $n$ of row $n$ — one cell per row, the cells boxed one at a
time, never a line drawn through them. The seven diagonal digits are
$\anchor{017.L.diagonal}$. Now the rule: the $n$th digit of the new
number $y$ is $1$ unless the diagonal digit is $1$, in which case it is
$2$ — the source says only ``$1$ or $2$, chosen so that it is different
from the $n$th digit of $x_n$'', and this is one way to choose. The
choices give $y = \anchor{017.L.y}$, both branches of the rule firing
(at positions one and six). Only now are the other digits of the rows
filled in, because they never mattered: $y$ was built from the diagonal
alone~\cite{probability-jason-filippou-a-better-number-sequence-to-present-cantor-s}.
Compare $y$ with row $n$ at column $n$ and they differ, row after row —
by construction, for every $n$, including the rows beyond the page,
since the rule is the same rule there. The picture is bookkeeping, not
geometry: any selection of one cell per row and at most one per column
would do the same
job~\cite{probability-jeremy-martin-cantor-s-diagonal-argument-math-410-kansas-200}.

\section{Why one or two}

Why not simply ``add one to the digit''? Because a rule that can write
$0$ or $9$ can land on a listed point under its other name. Take the
list $x_1 = 0.0999\ldots$, $x_2 = 0.999\ldots$, $x_3 = 0.999\ldots$,
and so on. Its diagonal is $0, 9, 9, \ldots$, and ``add one'' returns
$y = \anchor{017.N.plusone} = x_1$: different digits, the same number,
and the argument proves nothing~\cite{probability-michael-sipser-introduction-to-the-theory-of-computation-3rd}.
The binary flip fails the same way, since $0.0111\ldots_2$ and
$0.1000\ldots_2$ are both $1/2$; Cantor's own 1891 argument avoids the
issue by working on sequences of symbols rather than
numbers~\cite{probability-wikipedia-cantor-s-diagonal-argument}. With
its digits drawn from $\{1, 2\}$, $y$ ends in neither zeros nor nines,
so by the exception of the digits section it owns exactly one row.
That is the whole license: if $y$ were some $x_n$, the listed row of
$x_n$ would be a row for $y$ — its only row — so the two rows would
agree at digit $n$, where they were built to differ. A different digit
means a different number, \emph{because} $y$ never writes $0$ or $9$.
Two further comforts: $\anchor{017.C.range}$, so $y$ is inside $[0,1]$
(and both $1/2$ rows lie outside that range — the duplicate on the
list was harmless); and the objection ``just add $y$ to the list''
answers itself. Put $y$ at the front, shift the rows down, and run the
rule again: the new diagonal is $\anchor{017.M.rerun}$ and the new
number is $y' = \anchor{017.M.yprime}$ — different from $y$ at digit
one (always: $d'_1 = 3 - d_1$) and off every shifted row at its own
digit. The claim was never ``this list misses $y$''; it is ``every list
misses one''. Said directly, without the wrapper: for any sequence, the
rule produces a point not on
it~\cite{probability-andrej-bauer-proof-of-negation-and-proof-by-contradiction-20}.
Said as the source says it: ``the sequence $x_1, x_2, \ldots$ does not
exhaust the elements of $[0,1]$, contrary to what was assumed.''

\section{The same diagonal twice}

Cantor's 1891 paper never mentions
decimals~\cite{probability-cantor-ueber-eine-elementare-frage-der-mannigfaltigkeitslehr}.
Its rows are sequences of two symbols, $m$ and $w$ — $E^1 = (m, m, m,
\ldots)$, $E^2 = (w, w, w, \ldots)$, $E^3 = (m, w, m, w, \ldots)$ — and
the new row's $n$th symbol is whichever symbol row $n$'s $n$th symbol
is not. Two symbols, no representation to worry about; the reals came
afterwards, and the first uncountability proof (1874) used nested
intervals rather than a diagonal at
all~\cite{probability-georg-cantor-ueber-eine-eigenschaft-des-inbegriffs-crelle-77,probability-knapp-and-silva-the-uncountability-of-the-unit-interval-arxi}.
The same move works on subsets. List four subsets of $S = \{1, 2, 3,
4\}$ — Martin's example, $f(1) = \{1, 3\}$, $f(2) = \{1, 3, 4\}$, $f(3)
= \{\,\}$, $f(4) = \{2, 4\}$ — and draw the yes/no table whose entry
$(m, n)$ says whether $n$ is in the $m$th subset. Its diagonal reads
yes, no, no, yes. Flip it: the set $D$ that holds $n$ exactly when the
$n$th subset does not is $D = \anchor{017.Q.D}$, and it is on no row,
since it disagrees with row $n$ at $n$. Four subsets were listed and
$\anchor{017.Q.subsets}$ exist; no list of subsets of a set is ever
complete~\cite{probability-wikipedia-cantor-s-theorem}. The move recurs
once more in this repo's future: Turing ran it on the computable
sequences, where the fallacy lies in assuming the list itself can be
computed~\cite{probability-a-m-turing-on-computable-numbers-proc-lms-1936-37,probability-wikipedia-turing-s-proof}
— the undecidability story belongs to the algorithms topic, promised
there. And a smaller consequence, free: names are finite strings, and
finite strings are a sequence, so most points of $[0,1]$ have no name.

\section{Area, not sums}

Back to the wheel of fortune. What is the probability of a single
point? Suppose one point had $p = 0.15$; by symmetry so would every
other, and seven of them already stack to $\anchor{017.R.overflow} >
1$. Any $p > 0$ overflows with more than $1/p$ points, so a point has
probability zero — Example 1.4 of the
source~\cite{probability-bertsekas-and-tsitsiklis-introduction-to-probability-1st-ed}.
Now read the additivity axiom exactly as it is written: if $A_1, A_2,
\ldots$ is a \emph{sequence} of disjoint events, then $P(A_1 \cup A_2
\cup \cdots) = P(A_1) + P(A_2) + \cdots$. The finite additivity every
earlier chapter used is the same axiom's first clause; this is its
second, and it is stated for sequences — the same word Problem 4*
denies to $[0,1]$. A sequence of points therefore has probability $0 +
0 + \cdots = 0$, and the picture makes it visible without the axiom:
cover $x_1$ by an interval of length $\varepsilon/2$, $x_2$ by
$\varepsilon/4$, $x_3$ by $\varepsilon/8$, and the total length is
$\varepsilon$ — at $\varepsilon = 1/10$ the running totals are
$\anchor{017.R.cover}, \ldots$ — for any $\varepsilon$ at all. So if
$[0,1]$ were a sequence, $P([0,1])$ would be $0$, not $1$. The source
says so in a footnote: ``if the unit interval had a countable number of
elements, with each element having zero probability, the additivity
axiom would imply that the whole interval has zero probability, which
would contradict the normalization
axiom.''~\cite{probability-bertsekas-and-tsitsiklis-introduction-to-probability-1st-ed}
Probability is area on the interval \emph{because} the interval is not
a sum over points. Two guards close the chapter. This is consistency,
not existence: the argument shows a uniform law on $[0,1]$ is possible
only if the interval is uncountable, and that length \emph{is} a
probability law at all is the ``more advanced treatment'' the footnote
defers to — Lebesgue's, queued for the integration
series~\cite{probability-john-d-norton-the-material-theory-of-induction-2021-ch-14}.
And uncountable does not mean positive area: the Cantor set is
uncountable with length zero~\cite{probability-wikipedia-null-set}.
\end{narrative}

\section*{Practice problems}

\begin{problem}
Five points of $[0,1]$ are listed: $x_1 = 2/7$, $x_2 = 1/9$, $x_3 =
5/9$, $x_4 = 1/11$, $x_5 = 3/8$. Write each as a row of digits, read
the diagonal (digit $n$ of row $n$), and apply the chapter's rule ($1$
unless the diagonal digit is $1$, then $2$). Give the first five digits
of $y$ and check that $y$ differs from each row at its own digit.
\begin{hint}
$1/9 = 0.111\ldots$ and $1/11 = 0.0909\ldots$; the fourth digit of the
latter is a $9$, the fifth digit of $3/8 = 0.375$ is a $0$.
\end{hint}
\begin{solution}
The rows begin $0.28571$, $0.11111$, $0.55555$, $0.09090$, $0.37500$,
so the diagonal is $2, 1, 5, 9, 0$ and the rule gives $y = 0.12111
\ldots$ (the second digit is the branch to $2$). Column by column, $1
\ne 2$, $2 \ne 1$, $1 \ne 5$, $1 \ne 9$, $1 \ne 0$.
\end{solution}
\end{problem}

\begin{problem}
Put the $y$ of the previous problem at the front of its list (as $x_1$)
and shift the five rows down. Read the new diagonal — six digits now —
and the new number $y'$. Why must $y'$ differ from $y$ at the first
digit, whatever the list?
\begin{hint}
Only $y$'s first digit is consulted; the other five diagonal digits
come from the original rows, one column further right.
\end{hint}
\begin{solution}
The new diagonal is $y$'s first digit $1$, then $2/7$'s second digit
$8$, $1/9$'s third $1$, $5/9$'s fourth $5$, $1/11$'s fifth $0$, $3/8$'s
sixth $0$: $1, 8, 1, 5, 0, 0$, so $y' = 0.212111\ldots$. Its first
digit is $3 - d_1$: the rule sends $1$ to $2$ and anything else to $1$,
so applied to $y$'s own first digit, which is $1$ or $2$, it always
returns the other one.
\end{solution}
\end{problem}

\begin{problem}
Run the rule ``add one to the diagonal digit'' on the list $x_1 =
0.0999\ldots$, $x_2 = 0.999\ldots$, $x_3 = 0.999\ldots$, $\ldots$.
What number results, and why does this not show the list is
incomplete? Name the smallest change to the rule that repairs it.
\begin{hint}
Sum the geometric tails: $0.0999\ldots = (9/100)/(1 - 1/10)$.
\end{hint}
\begin{solution}
The diagonal is $0, 9, 9, \ldots$, so the rule returns $0.1000\ldots =
1/10$ — and $0.0999\ldots = (9/100)/(9/10) = 1/10$ is $x_1$ itself.
The new number differs from every row digit by digit and is still on
the list, because a row ending in nines has a second name. Choosing the
new digits from $\{1, 2\}$ (any two digits other than $0$ and $9$)
repairs it: the new number then has one name only.
\end{solution}
\end{problem}

\begin{problem}
Compute $0.111\ldots$ and $0.222\ldots$ as fractions, and show that
every number whose digits are all $1$s and $2$s lies between them.
\begin{hint}
Compare the two numbers digit by digit with $y$.
\end{hint}
\begin{solution}
$0.111\ldots = \tfrac{1/10}{1 - 1/10} = 1/9$ and $0.222\ldots = 2/9$.
For $y = \sum d_n 10^{-n}$ with every $d_n \in \{1, 2\}$, each term
satisfies $10^{-n} \le d_n 10^{-n} \le 2 \cdot 10^{-n}$, so summing,
$1/9 \le y \le 2/9$ — strictly inside $[0,1]$, whatever the list was.
\end{solution}
\end{problem}

\begin{problem}
Which points of $[0,1]$ have two decimal expansions? How many of them
have a terminating expansion of at most two digits after the point?
Write $1/2$ both ways.
\begin{hint}
Terminating decimals $k/10^m$ with $0 < k < 10^m$; count $m \le 2$ and
do not count the tenths twice.
\end{hint}
\begin{solution}
Exactly the terminating decimals strictly between $0$ and $1$: $k/10^m$
with $0 < k < 10^m$, written as $0.d_1 \cdots d_m 000\ldots$ or with
the last nonzero digit lowered by one and followed by nines. With at
most two digits they are $k/100$ for $k = 1, \ldots, 99$: ninety-nine
points, nine of which ($k$ a multiple of $10$) already have a one-digit
name. $1/2 = 0.5000\ldots = 0.4999\ldots$.
\end{solution}
\end{problem}

\begin{problem}
On the zigzag walk (anti-diagonals $p + q = 2, 3, \ldots$, each started
at the top row), which cell carries stamp $20$, and what fraction is
it? In the list of kept fractions (cells not in lowest terms skipped),
what position does $3/5$ take?
\begin{hint}
The stamp of $(p, q)$ is $(s-2)(s-1)/2 + p$ with $s = p + q$; list the
kept fractions in walk order and count.
\end{hint}
\begin{solution}
Stamp $20$ is cell $(5, 2)$, the fraction $5/2$ (its anti-diagonal
$s = 7$ starts at stamp $16$). The kept list runs $1/1, 1/2, 2/1, 1/3,
3/1, 1/4, 2/3, 3/2, 4/1, 1/5, 5/1, 1/6, 2/5, 3/4, 4/3, 5/2, 6/1, 1/7,
3/5, \ldots$, so $3/5$ (stamp $24$) is the nineteenth kept fraction.
\end{solution}
\end{problem}

\begin{problem}
Under the pairing of the first section (even $n \mapsto n/2$, odd $n
\mapsto -(n-1)/2$), which position does $-37$ occupy, and which integer
sits at position $100$?
\begin{hint}
Invert the two branches: a positive integer $z$ came from $2z$, a
non-positive one from $1 - 2z$.
\end{hint}
\begin{solution}
$-37$ is non-positive, so its position is $1 - 2(-37) = 75$ (check:
$75$ is odd and $-(75 - 1)/2 = -37$). Position $100$ is even and holds
$100/2 = 50$.
\end{solution}
\end{problem}

\begin{problem}
Let $S = \{1, 2, 3, 4, 5\}$ and list five subsets: $f(1) = \{2, 3\}$,
$f(2) = \{1, 2, 5\}$, $f(3) = \{\,\}$, $f(4) = S$, $f(5) = \{4\}$. Find
$D = \{s : s \notin f(s)\}$ and check it is none of the five. How many
subsets does $S$ have, and how many can any list of five reach?
\begin{hint}
Ask, for each $s$, whether $s$ lies in its own subset.
\end{hint}
\begin{solution}
$1 \notin f(1)$, $2 \in f(2)$, $3 \notin f(3)$, $4 \in f(4)$, $5 \notin
f(5)$, so $D = \{1, 3, 5\}$. It differs from $f(s)$ at $s$ for every
$s$: $1 \in D$ but not in $f(1)$, $2 \in f(2)$ but not in $D$, and so
on. $S$ has $2^5 = 32$ subsets; a list of five reaches at most five.
\end{solution}
\end{problem}

\begin{problem}[stretch]
Suppose every point of $[0,1]$ had probability $p = 0.03$. What is the
least number of points whose total exceeds $1$? Then cover a sequence
of points by intervals of lengths $\varepsilon/2, \varepsilon/4,
\varepsilon/8, \ldots$ with $\varepsilon = 1/100$: give the first four
running totals and say why the whole sequence has probability $0$.
\begin{hint}
$N > 1/p$; the running totals are $\varepsilon(1 - 2^{-k})$.
\end{hint}
\begin{solution}
$1/p = 33.3\ldots$, so $N = 34$: $34 \times 0.03 = 1.02 > 1$. The
totals are $1/200, 3/400, 7/800, 15/1600$, and every total stays below
$1/100$. The covered set has probability at most $\varepsilon$ for
every $\varepsilon > 0$, hence $0$ — which is what the additivity axiom
says directly: $0 + 0 + \cdots = 0$.
\end{solution}
\end{problem}
